\section{参数联动：同一个参数扫描究竟改变了什么}
本节起集中补足本轮精读所需的模型推导。前面已给完整两物种方程、广义欧姆定律、EMHD及四场陀螺流体，因此这里沿用那些定义，不重复列一套名字相近而归一化不同的方程。

\subsection{电子质量不能只保留在一条方程里}
式\eqref{eq:singlemom}中的相对漂移应力也可从动能直接理解。令$\vct w=\vel_i-\vel_e=\J/(en)$，把两物种速度代入其动能之和：
\begin{align}
 {n\over2}(m_iu_i^2+m_eu_e^2)
 &= {\rho u^2\over2}+{nm_im_e\over2M}w^2\\
 &= {\rho u^2\over2}+{m_im_e\over2Me^2n}J^2.
 \label{eq:selected-driftenergy}
\end{align}
交叉项恰为零，因为质心定义保证$\sum_s m_sn(\vel_s-\vel)=0$。
因此电流的惯性既储存能量，也携带动量通量。若保留电子尺度上的$\partial_t\J$，却任意丢掉配套对流项与应力，所得系统可能只是一种惯性正则化。
\citet{L160}对任意质量比的稳态分析保留了这类结构；其标量压力、对称稳态假设下，X点仍有$E_z=\eta J_z$，不能由“包含电子惯性”跳到“所有稳态率均与电阻无关”。

\subsection{温度、质量比、惯性尺度与FLR之间的约束}
本小节取单电离，$r_e=\sqrt{T_e/m_e}/|\Omega_{ce}|$、$\rho_s=\sqrt{T_e/m_i}/\Omega_{ci}$；这里$r_e$使用单分量热速度。
由定义逐步得到
\begin{equation}
 {r_e^2\over d_e^2}={\mu_0nT_e\over B_0^2}={\beta_e\over2},\qquad
 {r_e^2\over\rho_s^2}={m_e\over m_i},\qquad
 {\rho_s^2\over d_e^2}={\beta_e\over2}{m_i\over m_e}.
 \label{eq:selected-beta}
\end{equation}
所以固定$d_e,\rho_s$升高$\beta_e$，实际上必须降低$m_i/m_e$；固定$d_e,m_i/m_e$升高$\beta_e$，则$\rho_s$必增大。这正是\citet{E474}两类扫描不能直接比较的原因。
若另文用$v_{te}=\sqrt{2T_e/m_e}$定义回旋半径，则其$\rho_e=\sqrt2r_e$、$\rho_e/d_e=\sqrt{\beta_e}$。只看到一个$\sqrt2$差别，先查定义，不先判定物理矛盾。

在\citet{E474}的约定中，电子环平均因子为
\begin{equation}
 G_{10e}(k_\perp)=\exp(-k_\perp^2\beta_ed_e^2/4)
 =\exp(-k_\perp^2r_e^2/2).
\end{equation}
长波展开给$G_{10e}\simeq1+(r_e^2/2)\nabla_\perp^2$。这是从Fourier乘子独立获得的系数；所读预印本式(3.1)与这个定义存在因子2差别。有限$\beta_e$还伴随平行磁扰动及闭合改变，不能把全部变化归因于单一FLR项。

\section{最大熵高阶矩：从约束优化推导闭合}
\subsection{为什么再增加几个矩，并不能任意指定分布}
设我们只知道速度多项式$m_a(\vct v)$对应的矩
$U_a=\int m_af\,\dd^3v$。例如$1,v_i,v_iv_j$给十矩。选择在这些约束下使
\begin{equation}
 S[f]=-\int f(\ln f-1)\dd^3v
\end{equation}
最大的分布。这里省略常量参考分布以简化无量纲记号；它不改变变分步骤。
引入乘子$\alpha_a$，令
\begin{equation}
 \delta\left[S+\sum_a\alpha_a\left(\int m_af\dd^3v-U_a\right)\right]=0.
\end{equation}
因为$\delta S=-\int\ln f\,\delta f\dd^3v$，故
\begin{equation}
 f_\alpha=\exp(\alpha_am_a),\qquad
 \int m_af_\alpha\dd^3v=U_a.
 \label{eq:selected-maxent}
\end{equation}
指数形式保证$f>0$，但还要保证积分收敛。只有三次多项式作最高次时，指数一般沿某个方向趋于正无穷；因此保留热流而使用这类指数族时，需要适当的偶次最高项。14和21矩族增加$|\vct v|^4$约束，分别配合热流向量与完整三阶张量。\citep{E553}

十矩的二次指数可配方为各向异性Gaussian：
\begin{equation}
 f_{10}={n\over(2\pi)^{3/2}\sqrt{\det\Theta}}
 \exp\left[-{1\over2}(\vct v-\vel)^{\mathsf T}\Theta^{-1}(\vct v-\vel)\right],
 \qquad\Theta={\mathsf P\over mn}.
\end{equation}
$\Theta$须正定。这种分布允许不同方向温度不同，却不能产生两个分离峰；知道压力张量不等于知道所有粒子轨道群。

\subsection{如何求系数，以及为何具有对称双曲结构}
定义$h(\alpha)=\int f_\alpha\dd^3v$。以\emph{固定的目标矩}$U$最小化
\begin{equation}
 \mathcal L(\alpha;U)=h(\alpha)-\alpha\cdot U,\qquad
 \partial_{\alpha_a}\mathcal L=\int m_af_\alpha\dd^3v-U_a.
\end{equation}
不能在求导时又把$U$换成依赖$\alpha$的矩，否则改变了优化问题。
其Hessian为
\begin{equation}
 H_{ab}=\int m_am_bf_\alpha\dd^3v,\qquad
 \vct z^{\mathsf T}H\vct z=\int(z_am_a)^2f_\alpha\dd^3v>0
\end{equation}
（对线性独立基和可实现内部状态）。Newton步解
$H\delta\alpha=U-\int\vct mf_\alpha\dd^3v$，仍需速度空间积分和维持可积性的步长控制。

令$j_k(\alpha)=\int v_kf_\alpha\dd^3v$，矩守恒律可写为
\begin{equation}
 h_{\alpha\alpha}\partial_t\alpha+
 \sum_k(j_k)_{\alpha\alpha}\partial_{x_k}\alpha=S(\alpha).
\end{equation}
前矩阵正定，空间矩阵对称，因而在适当参数域内是对称双曲系统。这解释了数值上的优势；但非线性优化代价、可实现集合边界和高阶矩缺失仍在。
\citet{E553}只从PIC取矩后重建分布，不能据此宣称已求解上述21矩演化系统。

\subsection{跟算例：只约束到二阶时，最大熵分布为什么是高斯}
为把前面的抽象乘子变成可计算的例子，设已知数密度$n$、平均速度$\vel$和正定协方差
$\Theta_{ij}=P_{ij}/(mn)$。最大化$-\int f\ln(f/f_*)\dd^3v$，同时固定
$\int f$、$\int v_if$和$\int v_iv_jf$。对$f$变分后，驻值条件具有形式
\begin{equation}
 -\ln(f/f_*)-1+\lambda_0+\lambda_iv_i+\Lambda_{ij}v_iv_j=0.
\end{equation}
因此$f$是速度二次式的指数。为了使$\int f$有限，二次型$\Lambda$必须负定。把指数完成平方、用给定均值与协方差确定系数：
\begin{equation}
 f(\vct v)={n\over(2\pi)^{3/2}\sqrt{\det\Theta}}\,
 \exp\left[-{1\over2}(\vct v-\vel)^T\Theta^{-1}(\vct v-\vel)\right].
\end{equation}
若$\Theta=(T/m)\mathsf I$，便退回各向同性Maxwell分布。若$\Theta$有不同特征值，则沿不同主轴具有不同温度宽度。

该高斯关于$\vct c=\vct v-\vel$为偶函数，所以所有三阶中心矩
$Q_{ijk}=m\int c_ic_jc_kf\dd^3v$为零。这并不证明真实无碰撞重联的热流为零；它说明\emph{仅给二阶矩并选这一最大熵闭合}时，所重建分布没有偏斜和多束信息。若想保留非零三阶矩，就必须再加约束并处理更高阶可积性问题。

另一个实用结论是：若数值压力矩阵出现负特征值，$\Theta^{-1}$仍可能在代数上存在，但这份高斯已经不能作为物理概率密度。矩的“可实现性”就是要求存在非负、可积的分布产生这些矩，而不是仅要求程序能够对矩阵求逆。


\section{从双流体到KREHM：电子加热为何需要额外自由度}
\subsection{排序、未知量与离子响应}
\citet{E512}不是一般Hall MHD，而是强导场陀螺动理学的低$\beta_e$约化。关键序为
\begin{equation}
 \omega\ll\Omega_{ci},\quad k_\parallel/k_\perp\ll1,\quad
 \delta B_\perp/B_0\ll1,\quad \beta_e\sim m_e/m_i\ll1,
 \quad k_\perp d_e\sim k_\perp\rho_s\sim1.
\end{equation}
其中允许尺度之比作后续渐近展开；并不是要求每个算例$d_e$与$\rho_s$严格相等。取$T_i/T_e=\tau$为有限量，保留离子FLR；电子FLR则因$k_\perp\rho_e\ll1$舍去。下列式子改写成SI，且取$Z=1$。

记$a=A_\parallel$、$\delta\B_\perp=\nabla a\times\hat z$，定义
\begin{equation}
 D_t=\partial_t+{1\over B_0}[\phi,\cdot],\qquad
 \nabla_\parallel=\partial_z-{1\over B_0}[a,\cdot],\quad
 [f,g]=\partial_xf\partial_yg-\partial_yf\partial_xg.
\end{equation}
在这一排序下，离子非绝热分布$g_i$满足齐次平流。选择初值$g_i=0$便一直为零；这是一个子类，不是任意初值下离子动理学都消失。
两次环平均产生
\begin{equation}
 N\equiv{\delta n_e\over n_0}=-{1\over\tau}(1-\Gamma_0){e\phi\over T_0},\quad
 \Gamma_0(b)=I_0(b)e^{-b},\quad b={k_\perp^2\rho_i^2\over2},
 \label{eq:selected-quasi}
\end{equation}
这里$\rho_i=\sqrt{2T_i/m_i}/\Omega_{ci}$，$T_0=T_{e0}$。
长波$1-\Gamma_0\simeq b$，于是$N\simeq(e\rho_s^2/T_0)\nabla_\perp^2\phi$。
同时$u_{\parallel i}=0$至该阶，Ampère定律给
\begin{equation}
 u_{\parallel e}={\nabla_\perp^2a\over\mu_0en_0}
 ={e d_e^2\over m_e}\nabla_\perp^2a.
\end{equation}
特别冷的离子还需检查碰撞排序：原文要求$\beta_e\ll\tau^3$以忽略离子碰撞，不能无限制令$\tau\to0$。

\subsection{扣除密度和流速矩，剩下的就是缺失信息}
令$F_0=n_0(\pi v_{te}^2)^{-3/2}e^{-v^2/v_{te}^2}$，这里$v_{te}=\sqrt{2T_0/m_e}$。
把电子扰动写为
\begin{equation}
 \delta f_e=\left(N+{2v_\parallel u_{\parallel e}\over v_{te}^2}\right)F_0+g_e,
 \qquad\int g_e\dd^3v=\int v_\parallel g_e\dd^3v=0.
\end{equation}
因此$g_e$不重复计算密度和平均流；它储存温度、热流及更高矩。记
\begin{equation}
 \theta={\delta T_{\parallel e}\over T_0}
 ={1\over n_0}\int{2v_\parallel^2\over v_{te}^2}g_e\dd^3v.
\end{equation}
对约化电子动理学方程取零阶和一阶矩，再将它们从原方程减去，得
\begin{align}
 D_tN&=-\nabla_\parallel u_{\parallel e},\label{eq:selected-krehn}\\
 D_t(a-d_e^2\nabla_\perp^2a)&=D_\eta\nabla_\perp^2a-\partial_z\phi
 +{T_0\over e}\nabla_\parallel(N+\theta),\label{eq:selected-kreha}\\
 D_tg_e+v_\parallel\nabla_\parallel(g_e-\theta F_0)
 &=C[g_e]+\left(1-{2v_\parallel^2\over v_{te}^2}\right)
 F_0\nabla_\parallel u_{\parallel e}.\label{eq:selected-krehg}
\end{align}
对选用的简化碰撞算子，$D_\eta=\nu_{ei}d_e^2$。
读者可对最后一式积分检查两个零矩约束：右端源的零阶矩为零，一阶矩因奇偶性为零。
若直接令$g_e=0$，右端一般还含非零的$\nabla_\parallel u_{\parallel e}$；所以等温双流体通常是额外闭合，不是这组方程自动保持的解。

\subsection{温度梯度交换能量，但这里不独自断开磁通}
将式\eqref{eq:selected-kreha}改写为
\begin{equation}
 \partial_ta=-\nabla_\parallel\widetilde\phi
 +D_\eta\nabla_\perp^2a+D_t(d_e^2\nabla_\perp^2a),\qquad
 \widetilde\phi=\phi-{T_0\over e}(N+\theta).
\end{equation}
当最后两项忽略时，取旋度得到普通磁场随
$\vel_{\rm eff}=\hat z\times\nabla_\perp\widetilde\phi/B_0$冻结。
因此该约化中的电子温度项能改变输运速度和能量去向，却不是独立的破冻结项。
此结论依赖该低$\beta$形式，不能扩展到任意三维压力张量。

电子分布的二次自由能包含
\begin{equation}
 Y={n_0T_0\over2}\int N^2\dd V,\quad
 K_e={1\over2\mu_0}\int d_e^2(\nabla_\perp^2a)^2\dd V,\quad
 H_e=\int\!\dd V\dd^3v\,{T_0g_e^2\over2F_0}.
\end{equation}
另有磁能$U=\int|\nabla_\perp a|^2/(2\mu_0)\dd V$和离子极化能$K_i$。
$K_i$在长波冷离子极限回到$\int m_in_0u_E^2/2\dd V$；一般则须保留$1-\Gamma_0$的谱权重。
令$W_f=K_i+Y+U+K_e$。由矩方程乘对应场再分部积分，电子平行压力做功项在两部分之间相消，得到
\begin{align}
 {\dd W_f\over\dd t}&=-Q-\int\eta j_\parallel^2\dd V,\\
 {\dd H_e\over\dd t}&=Q-D_c,\quad
 Q=\int u_{\parallel e}\nabla_\parallel\delta p_{\parallel e}\dd V,\\
 D_c&=-\int\!\dd V\dd^3v\,{T_0g_eC[g_e]\over F_0}\ge0.
 \label{eq:selected-heating}
\end{align}
周期或无净通量边界是这里积分的前提。$Q$可正可负，是可逆交换；$D_c$和欧姆热才对应该模型中正定的热化通道。不能把$Y$增加直接叫温升。这里讨论连续光滑模型的预算；若另引入激波弱解、物理黏性或其他碰撞闭合，须把相应熵源另行列入。

\subsection{Hermite展开把相混合变成可计算的矩链}
对$v_\perp$积分，令$\xi=v_\parallel/v_{te}$，将剩余分布写成
\begin{equation}
 g_e=F_0\sum_{m=2}^{\infty}{H_m(\xi)\over\sqrt{2^m m!}}g_m,
 \quad H_0=1, H_1=2\xi, H_2=4\xi^2-2,\quad g_2={\theta\over\sqrt2}.
\end{equation}
利用$\xi H_m=H_{m+1}/2+mH_{m-1}$，流送$v_\parallel\nabla_\parallel g_e$只耦合相邻阶：
\begin{align}
 D_t\theta+\sqrt3v_{te}\nabla_\parallel g_3&=-2\nabla_\parallel u_{\parallel e},\\
 D_tg_m+v_{te}\nabla_\parallel\left(\sqrt{{m+1\over2}}g_{m+1}
 +\sqrt{{m\over2}}g_{m-1}\right)&=-\nu_{ei}m g_m,\quad m\ge3.
\end{align}
碰撞算子特意不阻尼温度矩$m=2$，使所需能量守恒。正交性给
\begin{equation}
 H_e={n_0T_0\over2}\int\sum_{m\ge2}g_m^2\dd V,\qquad
 D_c=n_0T_0\nu_{ei}\int\sum_{m\ge3}m g_m^2\dd V\ge0.
\end{equation}
由此可直接检验耗散符号。所读版本式(110)--(112)给$D_c$的负号与其式(83)、H定理不一致；本报告依上述定义使用正号，不据此推断作者代码有同样问题。

半碰撞$\nu_{ei}\gg\omega,|k_\parallel|v_{te}$时，$g_4\ll g_3$，从$m=3$方程解出
\begin{equation}
 g_3\simeq-{v_{te}\over2\sqrt3\nu_{ei}}\nabla_\parallel\theta,\qquad
 D_t\theta={v_{te}^2\over2\nu_{ei}}\nabla_\parallel^2\theta-2\nabla_\parallel u_{\parallel e}.
\end{equation}
于是热传导闭合是碰撞主导的渐近结果，而不是所有弱碰撞状态都可使用的公式。

\subsection{有限加热极限为何需要足够时间}
对线性沿场相混合的高$m$支，定义$E_m=|g_m|^2/2$，其大$m$连续近似为
\begin{equation}
 \partial_tE_m=-|k_\parallel|v_{te}\partial_m(\sqrt{2m}E_m)-2\nu_{ei}mE_m.
\end{equation}
稳态积分给
\begin{equation}
 E_m={C\over\sqrt m}\exp[-(m/m_c)^{3/2}],\qquad
 m_c=\left({3|k_\parallel|v_{te}\over2\sqrt2\nu_{ei}}\right)^{2/3}.
\end{equation}
耗散$2\nu_{ei}\int mE_m\dd m$正比$\nu_{ei}m_c^{3/2}$，因而在输入幅度$C$有有限极限时可不依赖$\nu_{ei}$。
但若$E_m$同时以$2\gamma$增长，解另含
\begin{equation}
 \exp[-(m/m_\gamma)^{1/2}],\qquad
 m_\gamma=\left({|k_\parallel|v_{te}\over2\sqrt2\gamma}\right)^2.
\end{equation}
当$m_\gamma\ll m_c$，分布还没把足够自由能送到碰撞尺度，就已随不稳定性增长；此时固定$\gamma$令$\nu\to0^+$，耗散趋零。
这不是“Landau阻尼消失”：场与分布间仍可交换能量，只是不可逆热化取决于极限顺序与观察时间。
数值上还必须令Hermite截断超过有效截止并检验闭合反射，否则人工截断会把相混合能量反送回低阶矩。

\subsection{补算：Hermite相邻矩耦合的平方根系数}
取$\xi=v_\parallel/v_{te}$，本小节$v_{te}=\sqrt{2T_e/m_e}$，与使用$\sqrt{T_e/m_e}$的记号相差$\sqrt2$。物理学家约定的Hermite多项式满足
$H_{m+1}=2\xi H_m-2mH_{m-1}$。令
$\varphi_m=H_m/\sqrt{2^m m!}$，逐项除以归一化因子：
\begin{equation}
 \xi\varphi_m=
 \sqrt{{m+1\over2}}\varphi_{m+1}
 +\sqrt{{m\over2}}\varphi_{m-1}.
\end{equation}
于是自由流动项$v_\parallel\nabla_\parallel g$投影后必然把第$m$矩与$m\pm1$矩联系起来。只展示这一项的矩链是
\begin{equation}
 \partial_tg_m+
 v_{te}\nabla_\parallel
 \left(\sqrt{{m+1\over2}}g_{m+1}
       +\sqrt{{m\over2}}g_{m-1}\right)=0.
\end{equation}
完整KREHM还含场耦合、低阶矩扣除、非线性对流和碰撞；这一行只是解释流动耦合系数的来源，不能单独当作完整重联模型。

\subsection{相混合可以先搬运信息，再在很小碰撞下产生热}
均匀场下一个沿场Fourier模的自由流动方程为
$\partial_tg_k+ik_\parallel v_\parallel g_k=0$，解为
\begin{equation}
 g_k(v_\parallel,t)=g_k(v_\parallel,0)
 e^{-ik_\parallel v_\parallel t}.
\end{equation}
随着时间增大，分布在速度空间的相位变化尺度变小：
$\Delta v_\parallel\sim1/(|k_\parallel|t)$。低阶速度积分可以因相位抵消而变小，但细结构的信息仍保留在分布中，因此场或低阶矩衰减不等于立刻不可逆热化。

若存在速度扩散$D_v\partial_{v_\parallel}^2g$，其有效作用率量级为
$D_v/\Delta v_\parallel^2\sim D_vk_\parallel^2t^2$。时间积分达到一量级时，
\begin{equation}
 \int_0^{t_c}D_vk_\parallel^2t^2\dd t
 \sim{D_vk_\parallel^2t_c^3\over3}\sim1,
 \qquad t_c\sim(D_vk_\parallel^2)^{-1/3}.
\end{equation}
这里$D_v$的单位是速度平方除以时间，是一个用于说明相混合的局部模型系数；并非把真实库仑碰撞算子全部化成常系数扩散。这个估计说明为什么很弱碰撞在足够长时间后仍会作用于很细的速度结构。

在Hermite语言里，这对应自由能从低$m$传到高$m$。有限矩截断若简单令最高层以上为零，会使矩空间的传输反射回来，产生人为回波。加入某种闭合是在选择怎样处理被截去的信息；需要验证它是否正确表示相混合与碰撞，而不只是让计算更稳定。


\section{线性双层与非线性爆发：两种不同的增长问题}
\subsection{无碰撞撕裂的两个内层尺度}
\citet{E512}附录B从外区理想方程到双层匹配提供了一条完整路线。取平衡$B_y/B_0\simeq x/L_s$，外区仍满足
$a_1''-[k_y^2+B_y''/B_y]a_1=0$，用$\Delta'=[a_1']/a_1(0)$连接内层。
定义电子共振尺度$\delta=\gamma L_s/(k_yv_{te})$、匹配尺度
$\delta_{\rm in}=(\sqrt2\rho_s\delta)^{1/2}$，再定义
\begin{equation}
 \lambda^2={\delta_{\rm in}^2\over d_e^2}={\gamma L_s\over k_y\VA d_e},\qquad
 \alpha={\rho_s\over d_e\lambda}.
\end{equation}
本小节$\VA$用导场$B_0$计算，$L_s$包含弱重联场的剪切，不可再额外乘一次重联场比。
当$\alpha\gg1$，匹配积分集中在电子层，量级关系
$\lambda^2\sim\Delta'\delta_{\rm in}\alpha$给
\begin{equation}
 \gamma\sim{k_y\VA\over L_s}d_e\rho_s\Delta',\quad
 \delta\sim d_e^2\Delta',\quad
 \delta_{\rm in}\sim d_e(\rho_s\Delta')^{1/2}.
\end{equation}
这一步可以逐项代入上面三个定义检查。常磁通近似要求$\Delta'\delta_{\rm in}\ll1$。
大$\Delta'$时以$1/\delta_{\rm in}$替代$\Delta'$的量级，得到
\begin{equation}
 \delta_{\rm in}\sim d_e^{2/3}\rho_s^{1/3},\quad
 \delta\sim d_e^{4/3}\rho_s^{-1/3},\quad
 \gamma\sim{k_y\VA\over L_s}d_e^{1/3}\rho_s^{2/3}.
\end{equation}
因此“内层宽度”至少可能指两个不同量，不能统一标成$d_e$。
原文完整匹配还给$\Delta'\delta_{\rm in}/2=I_G\Lambda/(1-I_G\Lambda^2)$，
其中$I_G$由非等温电子响应积分决定；等温时$I_G=\pi/2$，热离子处理另依赖Padé近似，不能把数值系数当成任意分布的精确常数。

半碰撞情形在这套线性内层方程中对应$d_e^2\mapsto D_\eta/\gamma$，并须同时改变温度响应函数。
这是一种方程结构映射，不是在完整非线性模拟中把质量换成电阻的通用规则。
\citet{L210}进一步把最快局地率换算为全局Alfvén时钟，讨论惯性“理想撕裂”；这里“理想”指增长达到宏观时间尺度，电子惯性并未被删掉。

\subsection{从势能下降推导爆发指数}
设$\epsilon$是描述岛扩张或流体位移的集体坐标，约化能量为
$H=M(\epsilon)\dot\epsilon^2/2+U(\epsilon)$。
变分得到
\begin{equation}
 M\ddot\epsilon+\tfrac12M'\dot\epsilon^2+U'=0.
\end{equation}
\citet{L191}在所构造的冷双场试探流族内，非线性阶段有效$M$近似常量、$U\propto-\epsilon^3$，所以
$\ddot\epsilon\propto\epsilon^2$，能量守恒又给$\dot\epsilon\propto\epsilon^{3/2}$。
这比线性指数增长快，却不同于早期试探速度场导出的$\ddot\epsilon\propto\epsilon^4$；差别涉及内层映射和边界能量项，不是简单拟合误差。

暖双场局域X结构中，\citet{L214}用局地位移$\sigma$及外部岛尺度$\epsilon$，得到
$\dot\sigma\propto\sigma^2$、$\sigma\propto\epsilon^{3/4}$。链式法则给
\begin{equation}
 \dot\epsilon={4\over3}{\epsilon\over\sigma}\dot\sigma
 \propto\epsilon\sigma\propto\epsilon^{7/4}.
\end{equation}
因此比较“爆发指数”必须先统一幅度定义。
对$\dot\epsilon=C\epsilon^p$，积分是
\begin{equation}
 \epsilon(t)=\left[\epsilon_0^{1-p}-(p-1)C(t-t_0)\right]^{-1/(p-1)}.
\end{equation}
形式有限时间发散只说明约化中间渐近解会离开适用区间；实际有限磁能、岛尺寸、次级磁岛和网格截止都会终止它。
\citet{F697}虽找到暖模型相近的$7/4$规律，也指出冷模型次级磁岛使增长间歇化，且早期较少横向模数漏掉了这些岛。

\section{把MRX的慢电子排出写成一个可检验的方程}
\citet{L119}采用的估计不是普适重联定律，而是稳态电子排出动量预算。
在忽略排出方向压力梯度及静电场、近似线性加速且密度平缓时，沿电子层长$L_e$积分得
\begin{equation}
 {m_en v_o^2\over2}\simeq{B_m^2\over2\mu_0}
 -{m_en\nu v_oL_e\over2}.
\end{equation}
磁张力项的简化用了入口/出口电场相等与连续性：$v_{
m in}B_m=v_oB_n$、$v_{
m in}L_e=v_o\delta_e$，所以$B_nL_e=B_m\delta_e$。
除以$B_m^2/(2\mu_0)$，令$u=v_o/V_{\!Ae}$、$A=\nu L_e/V_{\!Ae}$，便有
\begin{equation}
 u^2+Au=1,\qquad u={\sqrt{A^2+4}-A\over2}.
\end{equation}
$A\ll1$时$u\simeq1-A/2$，$A\gg1$时$u\simeq1/A$。
实验经典碰撞估计只能降到约$0.3$--$0.6V_{\!Ae}$；使用从$E/J$反推的有效阻尼后更接近观测$0.11V_{\!Ae}$，但这一步引入了经验输入及各方向阻尼相同的假设。
压力、静电场和涨落项未被独立排除，因此该计算解释了\emph{为何可能变慢}，尚未证明\emph{实际由哪一项主导}。
